\documentclass[11pt,a4paper]{article}

\usepackage[margin=2.2cm]{geometry}
\usepackage{booktabs}
\usepackage{amsmath}
\usepackage{listings}
\usepackage{xcolor}
\usepackage{hyperref}
\usepackage{longtable}
\usepackage{fancyhdr}

\definecolor{codebg}{RGB}{246,247,249}
\definecolor{codekw}{RGB}{20,101,95}
\definecolor{pass}{RGB}{20,101,95}
\definecolor{fail}{RGB}{156,69,38}

\lstdefinelanguage{GraphEasy}{
  keywords={graph, fn, void, int, while, if, print, kernel, numVertices, TRUE},
  keywordstyle=\color{codekw}\bfseries,
  sensitive=false,
  morecomment=[l]{//},
  morecomment=[s]{/*}{*/},
}

\lstset{
  language=GraphEasy,
  basicstyle=\ttfamily\small,
  backgroundcolor=\color{codebg},
  frame=single,
  framerule=0pt,
  breaklines=true,
  showstringspaces=false,
  tabsize=2,
  xleftmargin=8pt,
  xrightmargin=8pt,
  aboveskip=8pt,
  belowskip=8pt,
}

\newcommand{\passmark}{\textcolor{pass}{\textbf{PASS}}}
\newcommand{\cfg}[1]{\texttt{#1}}

\hypersetup{colorlinks=true, linkcolor=codekw, urlcolor=codekw}

\pagestyle{fancy}
\fancyhf{}
\rhead{\small DOACROSS PHI fix --- p1GraphEasy Autotuner}
\lhead{\small Test report}
\cfoot{\thepage}

\title{\textbf{DOACROSS / DOALL / Polly Trigger Suite}\\[0.3em]
\large Twenty cases, ten kernels, store-forwarded PHI fix}
\author{p1GraphEasy-con-AutoTuner \quad (\texttt{bench\_doacross\_phi\_fix.py})}
\date{2026-08-22 \quad N=48 vertices}

\begin{document}
\maketitle

\begin{abstract}
\noindent
We document the regression suite that validated the PDG fix for
\emph{store-forwarded scalar recurrence PHIs}: loops such as prefix sums
\texttt{a[i] = a[i-1] + \ldots} must classify as \textbf{DOACROSS}, not
\textbf{SEQUENTIAL}. The suite comprises \textbf{10} hand-written GraphEasy
kernels, each run under \textbf{two} compiler configurations
(\cfg{polly\_off} and \cfg{polly\_on}), for \textbf{20} compile-time checks.
All \textbf{20/20} passed. This report lists every kernel source, the
expected parallel class, measured IR markers, and the systems claim each case
supports.
\end{abstract}

\tableofcontents
\newpage

%==============================================================================
\section{What we are testing}
%==============================================================================

\subsection{Pipeline order (p1 Autotuner)}

\begin{verbatim}
IRGen -> canonicalize -> Polly O3 (+ -polly-parallel)
      -> Autotuner (layout ARS)
      -> PDG (my.loop.parallel = DOALL | DOACROSS | SEQUENTIAL)
      -> reconstructParallelIR -> loop outliner -> parallel_for_runtime
\end{verbatim}

Polly runs \emph{before} PDG. Affine nests may be consumed by Polly OpenMP
(\texttt{GOMP\_parallel*}); remaining loops are classified and optionally
outlined by the PDG path.

\subsection{The bug and the fix}

When LLVM mem2reg promotes a loop-carried scalar into a PHI, the PDG
dependence analysis may see a \emph{store-forwarded} pattern:
the load of \texttt{a[i-1]} is forwarded from the previous iteration's store.
Without special handling, such loops were marked \textbf{SEQUENTIAL} even though
they are classic \textbf{DOACROSS} (wavefront) candidates.

The fix in \texttt{pdg.cpp} recognizes these PHIs and assigns
\texttt{parallel.type=DOACROSS}. This suite proves the fix does not break
DOALL-only kernels and still composes with Polly on a combined program.

\subsection{Configurations}

\begin{center}
\begin{tabular}{ll}
\toprule
Config & Environment \\
\midrule
\cfg{polly\_off} & \texttt{GRAPH\_DISABLE\_POLLY=1} --- PDG/outliner own all loops \\
\cfg{polly\_on}  & Default Polly + \texttt{-polly-parallel} \\
\bottomrule
\end{tabular}
\end{center}

\subsection{How pass/fail is decided}

IR dumps at \texttt{DUMP\_LLVM\_BC\_PDG} and \texttt{DUMP\_LLVM\_BC\_AFTER\_OUTLINE}:
\begin{itemize}
  \item \textbf{DOACROSS expected:} \texttt{parallel.type=DOACROSS} $> 0$ in PDG dump,
        or \texttt{parallel\_for\_runtime(..., i32 1, i32 1)} after outline.
  \item \textbf{DOALL expected:} \texttt{parallel.type=DOALL} $> 0$ and
        \texttt{DOACROSS} $= 0$.
  \item \textbf{combo\_all3:} With Polly off: both DOALL and DOACROSS metadata.
        With Polly on: \texttt{GOMP} $> 0$ plus DOALL and DOACROSS still present
        on \emph{different} nests.
\end{itemize}

%==============================================================================
\section{Summary results (20/20 passed)}
%==============================================================================

\begin{center}
\small
\begin{longtable}{lllrrrrrr}
\toprule
Case & Config & Expect & OK & DOACROSS & DOALL & GOMP & pfr$_{11}$ & pfr$_{00}$ \\
\midrule
\endhead
01\_prefix\_pure & polly\_off & doacross & \passmark & 1 & 1 & 0 & 1 & 0 \\
01\_prefix\_pure & polly\_on  & doacross & \passmark & 1 & 1 & 0 & 0 & 0 \\
02\_prefix\_plus\_const & polly\_off & doacross & \passmark & 1 & 1 & 0 & 0 & 0 \\
02\_prefix\_plus\_const & polly\_on  & doacross & \passmark & 1 & 1 & 0 & 0 & 0 \\
03\_prefix\_plus\_x\_i & polly\_off & doacross & \passmark & 1 & 1 & 0 & 1 & 0 \\
03\_prefix\_plus\_x\_i & polly\_on  & doacross & \passmark & 1 & 1 & 0 & 0 & 0 \\
04\_prefix\_indirect & polly\_off & doacross & \passmark & 1 & 1 & 0 & 1 & 1 \\
04\_prefix\_indirect & polly\_on  & doacross & \passmark & 1 & 1 & 0 & 0 & 0 \\
05\_running\_mul & polly\_off & doacross & \passmark & 1 & 1 & 0 & 0 & 0 \\
05\_running\_mul & polly\_on  & doacross & \passmark & 1 & 1 & 0 & 0 & 0 \\
06\_saxpy\_doall & polly\_off & doall & \passmark & 0 & 1 & 0 & 0 & 1 \\
06\_saxpy\_doall & polly\_on  & doall & \passmark & 0 & 1 & 0 & 0 & 0 \\
07\_independent\_init & polly\_off & doall & \passmark & 0 & 1 & 0 & 0 & 0 \\
07\_independent\_init & polly\_on  & doall & \passmark & 0 & 1 & 0 & 0 & 0 \\
08\_fib\_like & polly\_off & doacross & \passmark & 1 & 1 & 0 & 1 & 0 \\
08\_fib\_like & polly\_on  & doacross & \passmark & 1 & 0 & 0 & 0 & 0 \\
09\_prefix\_then\_saxpy & polly\_off & doacross & \passmark & 1 & 1 & 0 & 0 & 0 \\
09\_prefix\_then\_saxpy & polly\_on  & doacross & \passmark & 1 & 1 & 0 & 0 & 0 \\
10\_triple & polly\_off & combo & \passmark & 1 & 1 & 0 & 1 & 3 \\
10\_triple & polly\_on  & combo & \passmark & 1 & 1 & 6 & 0 & 1 \\
\bottomrule
\end{longtable}
\end{center}

\noindent
\textbf{Legend:} \texttt{pfr$_{11}$} = outlined DOACROSS
(\texttt{parallel\_for\_runtime(..., 1, 1)}).
\texttt{pfr$_{00}$} = outlined DOALL (\texttt{..., 0, 0}).
\texttt{GOMP} = Polly \texttt{GOMP\_parallel*} sites in POST dump.

%==============================================================================
\section{What the suite proves (systems claims)}
%==============================================================================

\begin{enumerate}
  \item \textbf{PHI fix is sound.} Prefix-style recurrences (cases 01--05, 08--09)
        receive DOACROSS metadata with Polly off \emph{and} on; they are not
        misclassified as SEQUENTIAL.

  \item \textbf{No false DOACROSS on independent loops.} Saxpy and pure init
        (cases 06--07) stay DOALL-only (DOACROSS count zero).

  \item \textbf{Polly and PDG compose on different nests.} Case 10 with Polly on:
        six GOMP sites on the matmul nest, while DOALL and DOACROSS metadata
        remain on other loops --- not double-parallelization of the same nest.

  \item \textbf{Classification survives Polly transforms.} Even when Polly
        reshapes IR and outlining counts drop (\texttt{pfr$_{11}$}=0 with Polly on),
        PDG still attaches \texttt{parallel.type=DOACROSS} where expected.

  \item \textbf{Multi-operator programs.} Case 09 (prefix then saxpy) and case 10
        (matmul + saxpy + prefix) show mixed DOACROSS/DOALL regions in one
        \texttt{kernel()}.
\end{enumerate}

\noindent
\textbf{Not claimed by this suite:} end-to-end speedup. A separate ablation
(\texttt{bench\_parallel\_boost.py}) shows bare DOACROSS outlining can be
\emph{slower} than serial on large prefix loops; Polly wins primarily on dense
matmul-shaped nests.

%==============================================================================
\section{Per-case detail: code, expectation, proof}
%==============================================================================

Each listing is the \texttt{kernel()} body from the harness (graph header omitted:
\texttt{graph G \{ edges: file "edges.txt"; TRUE \};}).

% --- Case 01 ---
\subsection{Case 01: \texttt{01\_prefix\_pure}}
\textbf{Expect:} DOACROSS \quad \textbf{polly\_off:} \passmark \quad \textbf{polly\_on:} \passmark

\textbf{Proves:} Minimal inclusive prefix sum
\texttt{a[i] = a[i] + a[i-1]} is DOACROSS (canonical store-forwarded PHI).

\begin{lstlisting}
fn void kernel() {
  int n = numVertices(G);
  int a[n];
  int i = 0;
  while (i < n) { a[i] = i + 1; i = i + 1; }
  i = 1;
  while (i < n) { a[i] = a[i] + a[i - 1]; i = i + 1; }
  print a[n - 1];
}
kernel();
\end{lstlisting}

% --- Case 02 ---
\subsection{Case 02: \texttt{02\_prefix\_plus\_const}}
\textbf{Expect:} DOACROSS \quad \textbf{polly\_off:} \passmark \quad \textbf{polly\_on:} \passmark

\textbf{Proves:} Constant addend recurrence \texttt{a[i] = a[i-1] + 3} still DOACROSS
(not confused with independent update).

\begin{lstlisting}
fn void kernel() {
  int n = numVertices(G);
  int a[n];
  int i = 0;
  while (i < n) { a[i] = 1; i = i + 1; }
  i = 1;
  while (i < n) { a[i] = a[i - 1] + 3; i = i + 1; }
  print a[n - 1];
}
kernel();
\end{lstlisting}

% --- Case 03 ---
\subsection{Case 03: \texttt{03\_prefix\_plus\_x\_i}}
\textbf{Expect:} DOACROSS \quad \textbf{polly\_off:} \passmark \quad \textbf{polly\_on:} \passmark

\textbf{Proves:} Prefix with per-index operand \texttt{x[i]} --- dependence is still
on \texttt{a[i-1]}, not on cross-iteration \texttt{x}.

\begin{lstlisting}
fn void kernel() {
  int n = numVertices(G);
  int a[n];
  int x[n];
  int i = 0;
  while (i < n) { a[i] = 0; x[i] = i + 1; i = i + 1; }
  i = 1;
  while (i < n) { a[i] = a[i - 1] + x[i]; i = i + 1; }
  print a[n - 1];
}
kernel();
\end{lstlisting}

% --- Case 04 ---
\subsection{Case 04: \texttt{04\_prefix\_indirect}}
\textbf{Expect:} DOACROSS \quad \textbf{polly\_off:} \passmark \quad \textbf{polly\_on:} \passmark

\textbf{Proves:} Indirect read \texttt{x[Q[i]]} in the recurrence body does not
destroy DOACROSS classification of the \texttt{a[i-1]} chain.

\begin{lstlisting}
fn void kernel() {
  int n = numVertices(G);
  int a[n];
  int x[n];
  int Q[n];
  int i = 0;
  while (i < n) { a[i] = 1; x[i] = i + 1; Q[i] = n - 1 - i; i = i + 1; }
  i = 1;
  while (i < n) {
    int idx = Q[i];
    a[i] = a[i - 1] + x[idx];
    i = i + 1;
  }
  print a[n - 1];
}
kernel();
\end{lstlisting}

% --- Case 05 ---
\subsection{Case 05: \texttt{05\_running\_mul}}
\textbf{Expect:} DOACROSS \quad \textbf{polly\_off:} \passmark \quad \textbf{polly\_on:} \passmark

\textbf{Proves:} Multiplicative recurrence \texttt{a[i] = a[i-1] * 2} is DOACROSS
(same dependence shape as additive prefix, different operator).

\begin{lstlisting}
fn void kernel() {
  int n = numVertices(G);
  int a[n];
  int i = 0;
  while (i < n) { a[i] = 2; i = i + 1; }
  i = 1;
  while (i < n) { a[i] = a[i - 1] * 2; i = i + 1; }
  print a[n - 1];
}
kernel();
\end{lstlisting}

% --- Case 06 ---
\subsection{Case 06: \texttt{06\_saxpy\_doall}}
\textbf{Expect:} DOALL only \quad \textbf{polly\_off:} \passmark \quad \textbf{polly\_on:} \passmark

\textbf{Proves:} Independent \texttt{y[i] = 3*x[i] + y[i]} is DOALL with
\textbf{zero} DOACROSS (no false wavefront from the PHI fix).

\begin{lstlisting}
fn void kernel() {
  int n = numVertices(G);
  int x[n];
  int y[n];
  int i = 0;
  while (i < n) { x[i] = i + 1; y[i] = 1; i = i + 1; }
  i = 0;
  while (i < n) { y[i] = 3 * x[i] + y[i]; i = i + 1; }
  print y[n - 1];
}
kernel();
\end{lstlisting}

% --- Case 07 ---
\subsection{Case 07: \texttt{07\_independent\_init}}
\textbf{Expect:} DOALL only \quad \textbf{polly\_off:} \passmark \quad \textbf{polly\_on:} \passmark

\textbf{Proves:} Trivial parallel init \texttt{a[i] = i*2} --- baseline DOALL
with no recurrence.

\begin{lstlisting}
fn void kernel() {
  int n = numVertices(G);
  int a[n];
  int i = 0;
  while (i < n) { a[i] = i * 2; i = i + 1; }
  print a[n - 1];
}
kernel();
\end{lstlisting}

% --- Case 08 ---
\subsection{Case 08: \texttt{08\_fib\_like}}
\textbf{Expect:} DOACROSS \quad \textbf{polly\_off:} \passmark \quad \textbf{polly\_on:} \passmark

\textbf{Proves:} Second-order recurrence \texttt{a[i] = a[i-1] + a[i-2]} is
DOACROSS (two backward dependences, wavefront depth 2).

\begin{lstlisting}
fn void kernel() {
  int n = numVertices(G);
  int a[n];
  int i = 0;
  while (i < n) { a[i] = 0; i = i + 1; }
  a[0] = 1;
  if (n > 1) { a[1] = 1; }
  i = 2;
  while (i < n) { a[i] = a[i - 1] + a[i - 2]; i = i + 1; }
  print a[n - 1];
}
kernel();
\end{lstlisting}

% --- Case 09 ---
\subsection{Case 09: \texttt{09\_prefix\_then\_saxpy}}
\textbf{Expect:} DOACROSS \quad \textbf{polly\_off:} \passmark \quad \textbf{polly\_on:} \passmark

\textbf{Proves:} Sequential composition: DOACROSS prefix loop followed by DOALL
saxpy in one function --- PDG assigns correct classes per nest.

\begin{lstlisting}
fn void kernel() {
  int n = numVertices(G);
  int a[n];
  int y[n];
  int i = 0;
  while (i < n) { a[i] = 1; y[i] = 0; i = i + 1; }
  i = 1;
  while (i < n) { a[i] = a[i - 1] + 1; i = i + 1; }
  i = 0;
  while (i < n) { y[i] = a[i] + y[i]; i = i + 1; }
  print y[n - 1];
}
kernel();
\end{lstlisting}

% --- Case 10 ---
\subsection{Case 10: \texttt{10\_triple\_polly\_doall\_doacross}}
\textbf{Expect:} combo (Polly + DOALL + DOACROSS) \quad \textbf{polly\_off:} \passmark \quad \textbf{polly\_on:} \passmark

\textbf{Proves:} One program contains (1) dense matmul triple nest --- Polly GOMP
when on; (2) saxpy DOALL; (3) indirect prefix DOACROSS. Demonstrates that
Polly and PDG \emph{compose across nests}, not on the same nest.

\begin{lstlisting}
fn void kernel() {
  int n = numVertices(G);
  int A[n][n];
  int Bmat[n][n];
  int C[n][n];
  int x[n];
  int y[n];
  int pref[n];
  int Q[n];
  int i = 0;
  while (i < n) {
    int j = 0;
    while (j < n) {
      A[i][j] = i + 1;
      Bmat[i][j] = j + 1;
      C[i][j] = 0;
      j = j + 1;
    }
    x[i] = i + 1;
    y[i] = 1;
    pref[i] = 1;
    Q[i] = n - 1 - i;
    i = i + 1;
  }
  i = 0;
  while (i < n) {
    int j = 0;
    while (j < n) {
      int s = 0;
      int k = 0;
      while (k < n) {
        s = s + A[i][k] * Bmat[k][j];
        k = k + 1;
      }
      C[i][j] = s;
      j = j + 1;
    }
    i = i + 1;
  }
  i = 0;
  while (i < n) { y[i] = 3 * x[i] + y[i]; i = i + 1; }
  i = 1;
  while (i < n) {
    int idx = Q[i];
    pref[i] = pref[i - 1] + x[idx];
    i = i + 1;
  }
  print y[n - 1] + pref[n - 1] + C[1][1];
}
kernel();
\end{lstlisting}

\noindent
\textbf{Case 10 with Polly on:} \texttt{GOMP=6}, \texttt{DOALL=1}, \texttt{DOACROSS=1},
\texttt{pfr$_{00}$=1} --- matmul parallelized by Polly; PDG metadata still records
DOALL/DOACROSS on remaining loops.

%==============================================================================
\section{Reproduce}
%==============================================================================

\begin{verbatim}
cd p1GraphEasy-con-AutoTuner
./build_lowmem.sh
N=48 python3 bench_doacross_phi_fix.py
# -> benchmark_logs/doacross_fix_<timestamp>/REPORT.md
\end{verbatim}

\noindent
Related: trigger matrix (\texttt{bench\_polly\_vs\_pdg\_triggers.py},
\texttt{POLLY\_VS\_PDG\_TRIGGERS.md}); performance ablation
(\texttt{bench\_parallel\_boost.py}, \texttt{parallel\_boost\_REPORT.md}).

\end{document}
